\documentclass[11pt,a4paper]{article}
\usepackage{fontspec}
\usepackage{polyglossia}
\setdefaultlanguage{russian}
\setotherlanguage{english}
\setmainfont{Liberation Serif}
\setsansfont{Liberation Sans}
\setmonofont[Scale=0.9]{Liberation Mono}
\usepackage[a4paper,top=26mm,bottom=26mm,left=29mm,right=29mm]{geometry}
\usepackage{amsmath}
\usepackage{booktabs}
\usepackage{array}
\usepackage{longtable}
\usepackage{enumitem}
\usepackage{titlesec}
\usepackage{tikz}
\usetikzlibrary{arrows.meta,positioning,calc}
\usepackage[hidelinks]{hyperref}
\usepackage{microtype}

\setlength{\parindent}{0pt}
\setlength{\parskip}{5pt plus 1pt}
\linespread{1.04}
\setlength{\emergencystretch}{1.5em}

\titleformat{\section}{\normalfont\fontsize{12.5}{15}\bfseries}{\thesection.}{0.5em}{}
\titlespacing*{\section}{0pt}{14pt plus 2pt minus 2pt}{5pt}
\titleformat{\subsection}{\normalfont\fontsize{11}{13}\bfseries}{\thesubsection.}{0.5em}{}
\titlespacing*{\subsection}{0pt}{9pt plus 2pt minus 1pt}{3pt}

\setlist[enumerate]{leftmargin=2.2em,itemsep=1pt,topsep=3pt,parsep=0pt}
\setlist[itemize]{leftmargin=1.6em,itemsep=1pt,topsep=3pt,parsep=0pt,label=--}

\pagestyle{plain}

\tikzset{
  box/.style={draw,rectangle,thick,align=center,inner sep=6pt,font=\small},
  lbl/.style={font=\footnotesize,align=center},
  arr/.style={-{Latex[length=2.2mm]},thick},
  darr/.style={{Latex[length=2.2mm]}-{Latex[length=2.2mm]},thick},
}

\begin{document}

\begin{center}
{\fontsize{17}{21}\selectfont\bfseries Разработка программного обеспечения\\[2pt] роем автономных агентов\\[2pt] над единой средой исполнения\par}
\vspace{14pt}
{\normalsize Алексей Григорьев, Иван Опарин\\
Relux Works\\
\href{https://github.com/relux-works}{github.com/relux-works}\par}
\vspace{4pt}
{\normalsize Сентябрь 2026, третья редакция, версия 3.3\par}
\end{center}

\vspace{10pt}

\begin{list}{}{\leftmargin=1.2cm\rightmargin=1.2cm}
\item
\textbf{Аннотация.} Мы рассматриваем совместное изменение программного проекта агентами независимых владельцев в домене, где у принятия результата нет механического оракула: код, в отличие от формального доказательства, не проверяется ядром. Предмет работы есть разделение inference-ёмкости и полномочий на изменение проекта. Среда исполнения проекта существует в единственном авторитетном экземпляре, агентские циклы работают на машинах контрибьюторов под их собственными учётными данными и получают доступ только к санбоксу своей задачи. Рабочая нода может предлагать и производить изменения, но не может расширить область задачи, изменить условия её приёмки или продвинуть основную ветку. Авторизация, исполнение и принятие связаны проверяемыми записями реплицированного журнала, относящимися к конкретным версиям контракта задачи и состояния проекта: финальное одобрение и проверки привязаны к неизменяемому интеграционному кандидату, и основная ветка продвигается только к нему. Через консенсус проходят только детерминированные предикаты, суждение моделей остаётся снаружи, независимым и подотчётным. Протокол определён для трёх уровней доверия с одинаковыми инвариантами и разным определением кворума, и для каждого уровня сказано, от чего именно защищает сертификат. Мы приводим точные условия безопасности, механизмы, которые их обеспечивают, исполняемую модель фрагмента интеграции с исчерпывающим перебором состояний в заданных границах и план экспериментов, способных опровергнуть наши утверждения. Протокол не делает агентов умнее; он делает объединение чужой ёмкости на одном проекте безопасным и подотчётным.
\end{list}

\section{Введение}

Инструменты агентной разработки строятся вокруг одной схемы: агент работает там, где лежит код, изоляция достигается копиями рабочего дерева, параллелизм умножением инстансов. Для агентов это три проблемы. Среда бывает дорогой или невоспроизводимой: тулчейны, лицензии, ключи подписи, оборудование, данные, которым нельзя покидать машину. Много инстансов одного владельца упираются в лимиты провайдера модели. Агенты разных людей не могут работать на одной машине, не передавая ей свои учётные данные.

Из этого следует неравенство, которое мотивирует работу. Большие рои агентов, координируемые для работы над одним проектом, доступны сегодня вендорам моделей, владельцам крупных вычислительных ресурсов и корпорациям. Открытые проекты используют агентов, но устроены под координацию через людей. Мы предполагаем, что координация на уровне автономного роя даёт выигрыш в эффективности; это гипотеза, и раздел 18 описывает эксперимент, который её проверяет, а не утверждение.

Само объединение чужих агентов на общем проекте уже не является новым: в формальной математике платформы позволяют любому владельцу агента вносить доказательства в общую миссию, а совместная формализация крупных результатов на потребительских подписках показана на практике (раздел 19). В формальной математике у принятия есть оракул: ядро проверки доказательств. У кода его нет, и принятие приходится конструировать. Наш вклад есть конкретная композиция границ доверия вокруг совместного изменения состояния в домене без оракула: рабочая нода производит изменения, но не может расширить область задачи, изменить условия приёмки или продвинуть основную ветку; авторизация, исполнение и принятие связаны проверяемыми записями, относящимися к конкретным версиям контракта задачи и состояния проекта; суждение агентов вынесено за консенсус и сделано независимым и подотчётным.

Мы не утверждаем, что протокол улучшает суждение моделей: качество декомпозиции, вероятность обнаружения дефекта при ревью и сила проверок остаются свойствами моделей и проекта. Работа устроена так: точные условия безопасности (разделы 2, 4, 15), механизмы, которые их обеспечивают (разделы 5--14), исполняемая модель ключевого фрагмента (раздел 17) и эксперименты, способные нас опровергнуть (раздел 18).

\subsection*{Где этот документ сейчас}

Состояние на 29 сентября 2026 года. Протокол роя, описанный в этой работе, есть распределённый слой Apiary (пасеки). Пасека есть все узлы и процессы одного владельца; процессы разных владельцев работают вместе через внешних участников под арендами и подписями. Слой роя отвечает за этот случай: проект ведут агенты независимых владельцев над единой средой исполнения. С 24 сентября 2026 года большинство понятий этой работы получили собственные спецификации в репозиториях модулей. Где эта работа и спецификация расходятся, права спецификация. Работа сохраняет протокол, гарантии и модель: инварианты, уровни доверия с определением кворума, правила валидаторов, модель угроз, оценки и исполняемую модель фрагмента интеграции.

Спецификации лежат в организации relux-works и пока остаются черновиками: curator-playbook, \texttt{spec/process-configuration.md}, v4.8 (конфигурация процесса); curator-trust, \texttt{spec/trust.md}, v0.4 (люди, подтверждения, права); curator-keeper, \texttt{spec/keeper.md}, v0.5.1 (главный цикл узла); agent-session-host, \texttt{spec/session-host.md}, v0.3 (хост сессий); waggle, \texttt{spec/waggle.md}, v5.2 (подписанные сообщения, аренды, кворум). Борд и ядро процесса есть task-board в репозитории skill-project-management, выпуск v0.28.0. Решения записаны в вики экосистемы (\texttt{decisions.md}), порядок работ задаёт её роадмап (\texttt{roadmap/ecosystem-roadmap.ru.md}). Владельцы по разделам:

{\footnotesize
\begin{longtable}{p{3.4cm} p{10.0cm}}
\toprule
Раздел работы & Владелец сейчас; что остаётся здесь \\
\midrule
\endfirsthead
\toprule
Раздел работы & Владелец сейчас; что остаётся здесь \\
\midrule
\endhead
\bottomrule
\endfoot
\bottomrule
\endlastfoot
1. Введение & вики экосистемы: \texttt{README.md}, роадмап (черновик v4); здесь постановка задачи и вклад \\
2. Термины и уровни доверия & curator-trust v0.4 §1, §3, §13 (люди, ключи с присутствием человека, стадии L0--L2); имена для людей: вики, \texttt{naming.ru.md}; здесь кворум кор-нод на уровнях L0--L2 \\
3. Топология & curator-keeper v0.5.1 §3, §4 (узел, главный цикл, несколько узлов); здесь три роли узлов протокола \\
4. Инварианты & здесь, I1--I10; принципы ядра curator-playbook v4.8 §1 повторяют I2, I3, I4, I6, I7 \\
5. Среда исполнения & только здесь: санбокс, слоты, receipts, вход в санбокс; смежно curator-trust v0.4 §12 (отдельная учётная запись ОС для агентов) \\
6. Журнал & task-board v0.28.0 (потоки активности элементов); waggle v5.2 §4, §19 (подписанные конверты, акты на борде); здесь QC, дерево Меркла, секвенсер, светлый клиент \\
7. Контракт задачи и полномочия & curator-playbook v4.8 §2.4, §4.2, §5, §6 (\texttt{Playbook.json}, версия в силе, делегирование, охранные условия, наборы проверок); здесь digest плана и адрес аттестации \\
8. Правила валидаторов & curator-playbook v4.8 §6.3, §6.10, §10 (охранные условия ядра, инварианты доставки, статический валидатор); здесь предикаты перед подписью кор-ноды \\
9. Оркестратор & curator-keeper v0.5.1 §3.2, §4 (один корневой оркестратор на процесс, арендованные поддеревья); curator-playbook v4.8 §4.4 (корневая роль); task-board (запуск агентов, журнал целей); waggle v5.2 §12 (аренды областей) и трек CM2 (принимающий писатель); здесь аренда и смена эпохи \\
10. Цикл задачи & task-board (цикл продьюсер--ревьювер, Change Request, интеграция через compare-and-set); curator-playbook v4.8 §6.10, §6.12, §6.16; здесь кандидат $I$ и выбор ревьювера \\
11. Внешние эффекты & curator-playbook v4.8 §6.11--6.15 (исполнитель ядра, чеки, ключ идемпотентности, публикация); curator-trust v0.4 §18 (привилегированные операции) \\
12. Рабочие ноды & waggle v5.2 §3, §5; curator-trust v0.4; curator-keeper v0.5.1 §7.7 (существующие харнесы); agent-session-host v0.3; здесь ёмкость как ресурс роя \\
13. Независимость и репутация & waggle v5.2 §3, §4.5 (реестр ключей, уровни гарантии ключа); curator-trust v0.4 §3 (люди, сеть доверия ключей); здесь три оси, доказательства идентичности, репутация \\
14. Остановка и федерация & curator-keeper v0.5.1 §10.3 и agent-session-host v0.3 §9 (остановка узла); curator-trust v0.4 §5, §13; waggle v5.2, этап CM4 (федерация); здесь остановка среды и манифест проекта \\
15. Модель угроз & здесь; частично curator-trust v0.4 (ограждение и граница) и waggle v5.2 §7 (защита от инъекций) \\
16. Свойства и оценки & здесь; измерение $d$: curator-playbook v4.8 §6.16 \\
17. Проверяемый фрагмент & только здесь, модель в \texttt{paper/model/} \\
18--20. Эксперименты, смежные работы, заключение & только здесь \\
Приложение & здесь; бюджеты процесса: curator-playbook v4.8 §6.14 \\
\end{longtable}
}

Только в этой работе описаны вход в санбокс и среда исполнения, журнал с сертификатами кворума, секвенсером и светлым клиентом, исполняемая модель фрагмента интеграции и план экспериментов; других источников у этих частей нет.

\section{Термины и уровни доверия}

\textbf{Среда исполнения.} Авторитетная машина проекта: рабочее дерево, санбоксы, проверки, основная ветка. Может иметь исполнительные реплики и верификатор.

\textbf{Кор-ноды.} $n$ узлов с фиксированными ключами, принадлежащих независимым владельцам. Ведут журнал и проверяют правила.

\textbf{Секвенсер.} Сервис, упорядочивающий записи журнала. На нижнем уровне совпадает с единственной кор-нодой.

\textbf{Рабочие ноды.} Агентские циклы контрибьюторов на их машинах, с их inference-ёмкостью и ключами. Не доверены.

\textbf{Оркестратор.} Прикладная роль планирования, удерживаемая кор-нодой по аренде.

\textbf{Журнал, запись, QC.} Журнал есть упорядоченная цепочка подписанных записей. Запись совершена, когда собран сертификат кворума (QC) уровня.

\textbf{Эпоха.} Срок одной аренды оркестратора. Runtime агента и менеджер сессий используют это слово для времени жизни плана сессии и аренды владельца сессии; здесь эпоха всегда означает срок оркестратора.

\textbf{План.} Скомпилированный набор полномочий агентской сессии; digest плана фиксирует, что агенту разрешено.

\textbf{Контракт задачи.} Версионируемое описание того, что задача должна получить и по каким правилам результат принимается; digest контракта фиксирует условия приёмки (раздел 7; владелец определения сейчас есть конфигурация процесса).

\textbf{Кандидат.} Интеграционный кандидат $I$: неизменяемое состояние, построенное из ревизии $R$ над головой основной ветки $M$; предмет финального одобрения, проверок и интеграции.

\textbf{Аттестация.} Подписанный результат проверки свойства для конкретного кандидата, контракта, политики и окружения; исход есть пройдена, провалена или не установлена.

\textbf{Санбокс, слот, receipt, write scope.} Санбокс есть изолированный каталог и политика ограничения для одной задачи; слот есть право на одновременное исполнение; receipt есть подписанный факт исполнения команды; write scope есть множество путей, которые задача вправе изменять.

\textbf{Ключ идентичности, доказательство идентичности.} Долгоживущий ключ владельца ноды, привязанный к внешней идентичности (домен, аккаунт, запись в системе имён) двусторонним доказательством: подписанным заявлением владельца и артефактом у провайдера, содержащим отпечаток ключа. Ключи нод сертифицируются ключом идентичности и могут ротироваться, не меняя доказательств (раздел 13).

\textbf{Имена для людей.} Материалы для людей называют части системы именами пасеки, которые операторы утвердили 25 сентября 2026 года (вики, \texttt{naming.ru.md}); спецификации и эта работа пишут технические термины. Карта ставит рядом имя для людей, технический термин и то, что ему соответствует здесь. Плейбук и waggle сохраняют свои имена. Слово рой в этой работе описательное: так называется распределённый слой целиком, именем для людей оно не служит.

{\footnotesize
\begin{longtable}{p{2.1cm} p{3.1cm} p{8.0cm}}
\toprule
Имя для людей & Технический термин & В этой работе \\
\midrule
\endfirsthead
\toprule
Имя для людей & Технический термин & В этой работе \\
\midrule
\endhead
\bottomrule
\endfoot
\bottomrule
\endlastfoot
пасека & вся система & Apiary: все узлы и процессы одного владельца; протокол работы есть её распределённый слой \\
точка & узел & одна машина; на ней работает кор-нода, рабочая нода или среда исполнения \\
Кипер & главный цикл & прямого аналога нет; главный цикл узла запускает корневого оркестратора процесса \\
улей & процесс & проект со своей средой, журналом и оркестратором \\
матка & корневой оркестратор & оркестратор (раздел 9) \\
соты & борда & борд задач, в котором материализуется журнал (раздел 6) \\
пчела & запуск & агентская сессия рабочей ноды по одному назначению \\
дикие пчёлы & внешние участники & рабочие ноды контрибьюторов (раздел 12) \\
нуклеус & пакет процесса & прямого аналога нет; опубликованный плейбук, с которого начинается новый процесс \\
\end{longtable}
}

\textbf{Уровни доверия.} Инварианты раздела 4 одинаковы на всех уровнях. Различаются определение QC, порядок записей и то, от чего защищает сертификат.

\begin{center}
\footnotesize
\begin{tabular}{p{1.1cm} p{3.4cm} p{3.0cm} p{5.6cm}}
\toprule
Уровень & Кор-ноды и QC & Порядок & От чего защищает сертификат \\
\midrule
L0 & $n = 1$; QC есть подпись одного ключа & единственная нода & от действий вне кора; кор и среда доверены \\
L1 & $n \geq 2$; QC есть $t$ подписей, $t > n/2$ & доверенный секвенсер, валидаторы подписывают асинхронно & от одиночного несанкционированного решения при компрометации меньшинства ключей; от раздвоения истории защищает только честность секвенсера \\
L2 & $n \geq 3f+1$; QC есть $2f+1$ подписей & византийски устойчивый консенсус [1, 2, 4, 5] & и от несанкционированного решения, и от раздвоения истории при $f$ злонамеренных кор-нодах \\
\bottomrule
\end{tabular}
\end{center}

Порог большинства сам по себе не даёт единственности истории: при трёх подписантах и пороге два злонамеренный подписант может подписать две несовместимые записи, а два честных по одной каждый, и обе записи получат сертификат. Пересечение двух кворумов размера $t$ содержит не менее $2t - n$ участников, и чтобы в нём гарантированно был честный при $f$ злонамеренных, требуется $2t - n > f$. L1 этого не требует и потому доверяет секвенсеру в отношении порядка; L2 использует конкретный протокол консенсуса. Проект объявляет уровень в манифесте, повышение уровня есть запись о составе.

\textbf{Уровни и стадии доверия.} curator-trust (\texttt{spec/trust.md} v0.4, §13) называет стадиями L0--L2 другую ось: не кворум кор-нод, а то, кто подтверждает изменения процесса. Стадия L0 там есть один оператор и один авторитетный борд, L1 несколько операторов с подтверждениями M из N, L2 допуск внешних участников. L0 и L1 двух документов соответствуют друг другу, хотя порог M из N не обязан быть большинством. Внешние участники здесь суть рабочие ноды и допускаются на любом уровне, а византийскому консенсусу кор-нод на L2 в curator-trust соответствия пока нет. Подтверждение там требует ключа, которым нельзя подписать без присутствия человека: FIDO2 с проверкой пользователя, ключ в Secure Enclave, TPM или Windows Hello за аутентификацией ОС, приложение-компаньон. Граница проходит по присутствию человека, а не по проверкам внутри учётной записи ОС, от которой работают агенты. Для нескольких операторов на одном борде решено направление: один принимающий писатель с fencing по эпохе (раздел 9).

\section{Топология}

Система состоит из узлов трёх типов. Среда не запускает агентских циклов, она пассивна и исполняет команды, поступающие извне, только в границах, заданных журналом. Кор-ноды ведут журнал и выделяют оркестратора. Рабочие ноды выполняют задачи через единственный инструмент: удалённое исполнение команд в санбоксе задачи.

\begin{center}
\begin{tikzpicture}[node distance=10mm]
\node[box,minimum width=46mm,minimum height=16mm] (core) {Кор-ноды\\[1pt]\footnotesize журнал, правила, секвенсер, оркестратор};
\node[box,minimum width=46mm,minimum height=16mm,right=39mm of core] (work) {Рабочие ноды\\[1pt]\footnotesize агентские циклы контрибьюторов,\\собственная ёмкость и ключи};
\node[box,minimum width=64mm,minimum height=16mm,below=14mm of $(core.south)!0.5!(work.south)$] (env) {Среда исполнения\\[1pt]\footnotesize основная ветка, санбоксы, слоты, проверки, раннер операций};
\draw[darr] (core) -- node[lbl,above] {назначения, ревизии,\\вердикты, heartbeat} (work);
\draw[arr] (core.south) -- node[lbl,right,xshift=1mm] {записи с QC} (env.north -| core.south);
\draw[arr] (work.south) -- node[lbl,left,xshift=-1mm] {удалённое исполнение\\в санбоксе задачи} (env.north -| work.south);
\end{tikzpicture}
\end{center}

Каждая нода объявляет административный домен, семейство модели и тулчейн; какие из этих сведений подтверждаются, а какие принимаются на доверии, определяет политика независимости (раздел 13). Через сеть проходят только команды, их вывод, хеши и подписи. Учётные данные контрибьютора не покидают его машину, среда не покидает машину проекта.

\section{Инварианты}

Всё, что следует ниже, есть способы обеспечить десять инвариантов. Спецификации нижних уровней ссылаются на них по номерам.

\begin{enumerate}[label=I\arabic*.,leftmargin=2.6em]
\item Ни одного эффекта на среде без совершённой записи авторизации, digest плана которой совпадает с digest плана вызова, и актуальность которой проверена на границе исполнения.
\item Полномочия только сужаются: план задачи есть подмножество capability, capability есть часть цели, ревизия не изменяет ничего вне write scope; условия приёмки задачи меняются только новой версией контракта с отдельной авторизацией.
\item Основная ветка продвигается только атомарной заменой головы $M$ на неизменяемого кандидата $I$, к которому привязаны финальное одобрение и аттестации, ссылающиеся на текущую версию контракта и политики.
\item Суждение не входит в консенсус: валидаторы проверяют только детерминированные предикаты над совершёнными записями и хешами содержимого.
\item Продьюсер, ревьювер и исполнитель проверок одной задачи независимы согласно объявленной политике; отсутствие подходящих участников не ослабляет политику автоматически; случайность выбора фиксируется после заморозки кандидатов и хеша ревизии.
\item Ключи контрибьютора не покидают его машину, и подпись выполняется на машине агентского цикла; секреты проекта существуют только в раннере и недоступны ни из одного санбокса.
\item Журнал есть единственное состояние авторитета, сессия модели есть кэш; записи с устаревшей арендой отвергаются и сохраняются как расхождение.
\item Любая кор-нода останавливает систему, кворум уровня возобновляет; среда хранит наивысшую принятую версию журнала и не принимает более раннюю.
\item Ограничение есть сообщаемый факт: среда объявляет, что она способна ограничить, и задача, требующая полного ограничения, отказывает на частичном.
\item Все обращения к среде через протокол и выданные ею результаты имеют проверяемое происхождение и связаны цепочкой receipts; протокол не удостоверяет полноту внутреннего контекста недоверенной рабочей ноды.
\end{enumerate}

\section{Среда исполнения}

\textbf{Авторитетная среда.} Логически состояние одно: основная ветка и журнал определяют его. Физически среда может быть пулом исполнительных реплик, разделяющих рабочее дерево по снимкам; голова продвигается в одном месте. Верификатор есть реплика, независимо повторяющая критические проверки.

\textbf{Санбокс и слоты.} Санбокс задачи есть копия рабочего дерева на момент назначения, созданная копированием при записи, и политика ограничения на уровне операционной системы: запись только в каталог санбокса, чтение системы и тулчейна, запрет на остальной домашний каталог и сервисы, запрет сигналов и сведений о чужих процессах, сеть закрыта, кроме локального egress-прокси со списком разрешённых адресов. Копия рабочего дерева не является границей; граница есть механизм системы: Seatbelt на macOS, пространства имён, seccomp и cgroups на Linux, microVM там, где нужна прочность гипервизора. Агент никогда не получает shell хоста. Санбоксов может быть тысячи; одновременное исполнение ограничено слотами, и очередь есть часть протокола.

\textbf{Защищённые пути и раннер.} Проверки, конфигурация сборки, правила и манифест смонтированы во все санбоксы только на чтение; их изменение проходит через задачи привилегированного класса. Операции с секретами проекта и общими сервисами (подпись, публикация, устройства, данные) недоступны из shell и предоставляются как именованные операции раннера (раздел 11).

\textbf{Receipts и воспроизводимость.} Каждая команда получает идентификатор, digest команды и входного состояния, лимиты; по завершении среда выдаёт receipt с кодом выхода, дайджестами вывода, ресурсами и digest результирующего состояния. Receipts задачи образуют цепочку, голова которой входит в запись о ревизии (I10). Receipt удостоверяет происхождение, а не воспроизводимость. Воспроизводимость требует отдельно: сохранённых входов, digest окружения и тулчейна, зафиксированных зависимостей и, там, где они влияют на результат, внешних ответов; egress-прокси есть точка, где внешние ответы фиксируются. Среда публикует инвентарь ограничений; задача объявляет требуемый класс, несовпадение есть отказ (I9).

\textbf{Светлый клиент.} Среда исполняет только подтверждённое журналом: открывает санбокс по назначению с QC, проводит операцию раннера по её авторизации, продвигает голову по записи об интеграции с QC (I1). Она проверяет включение записи по доказательству Меркла [10] под подписанным снапшотом и хранит наивысшую принятую версию журнала как долговременное состояние безопасности (I8). Актуальность разрешения проверяется на границе исполнения: запись, совершённая в прошлой эпохе или до остановки, не даёт права на исполнение сейчас.

\textbf{Вход в санбокс.} Право на вход возникает только с записью назначения. При её совершении демон среды выпускает короткоживущий SSH-сертификат для транспортного ключа ноды: principal равен идентификатору задачи, срок равен окну назначения, принудительная команда запускает демон по стандартным потокам ровно для санбокса этой задачи, все прочие возможности SSH (перенаправление портов, agent forwarding, sftp, pty, переменные окружения) запрещены. Команда исполняется от непривилегированного служебного пользователя без секретов; отзыв не нужен, сертификат истекает сам, а закрытие задачи или остановка рвут сессии. Транспортный ключ сертифицируется ключом ноды и не совпадает с ключом подписи ревизий. Для публичного роя SSH-сервер встраивается в сам непривилегированный демон и принимает только исполнение протокола, так что sshd хоста и root в приёме соединений не участвуют. SSH здесь только аутентификация, связывающая сессию с назначением; границей остаётся санбокс, и вход безопасен ровно настолько, насколько безопасен он.

\textbf{Shell.} Для модели shell среды неотличим от локального: состояние сессии сохраняется между вызовами, каждый вызов есть отдельный процесс, мутирующие вызовы сериализованы, параллельность даётся именованными сессиями, длительные процессы живут по дескриптору, большой вывод остаётся в санбоксе. Это свойство качества, а не безопасности.

\section{Журнал}

Запись содержит хеш предыдущей, тип, тело или хеши тел и подписи; для записи $i$ с канонической формой $c(r_i)$ хеш есть $h_i = H(h_{i-1} \,\|\, c(r_i))$. Над хешами строится дерево Меркла, снапшот подписан и содержит корень, длину, голову и версию [11, 12]. История только дописывается. Порядок задаёт секвенсер, запись совершена при QC уровня, валидаторы подписывают только после проверки правил раздела 8. На L2 порядок фиксирует консенсус валидаторов; в любом случае среда исполняет только записи с QC. Лидер консенсуса есть механическая роль и не совпадает с оркестратором.

\textbf{Типы записей.} Цель, capability, контракт задачи (версии), назначение, зависимость (какие принятые артефакты задача использует и какие результаты предоставляет), ревизия, заморозка кандидатов в ревьюверы, случайность, кандидат, ревью кандидата, аттестация, интеграция, вестинг, аренда и эпоха, чекпоинт, операция раннера и её исход, пауза, возобновление, состав, аттестация репутации, подтверждение идентичности (свидетельство резолвера и его хеш), сертификация ключа ноды ключом идентичности.

\textbf{Доступность и материализация.} Тела хранятся адресуемо по содержимому, журнал содержит хеши, и хеш включается только с сертификатом доступности от независимых реплик. Код адресуется хешами коммитов, git выступает хранилищем. Граф зависимостей задач и артефактов есть представление журнала, материализуемое из записей о зависимостях, а не второй источник истины (I7); индекс для поиска годных к переиспользованию результатов помогает найти, но не определяет, принят ли результат и применим ли к текущей версии контракта. Журнал может материализоваться как борд задач в git-репозитории проекта: каждый переход есть коммит борда, QC покрывает его хеш. Записи в такой борд принимает один писатель с fencing по эпохе (раздел 9).

\section{Контракт задачи и полномочия}

Полномочия и условия приёмки есть два разных вопроса, и они фиксируются двумя разными объектами.

Владелец определений этого раздела сейчас есть конфигурация процесса (curator-playbook, \texttt{spec/process-configuration.md} v4.8). Проект объявляет в \texttt{Playbook.json} типы элементов и их workflow, охранные условия переходов, наборы проверок и делегирование с потолками; изменение процесса вступает в силу по подтверждению человека или в пределах выданных им прав (§2.4 спецификации). Где этот раздел и спецификация расходятся, права спецификация. Здесь остаётся протокольная форма: digest плана, который сверяет среда, и адрес аттестации из четырёх частей.

\textbf{Digest плана} отвечает на вопрос, что агенту разрешено: привязки инструментов, гранты, лимиты. Он компилируется из write scope и класса задачи, входит в запись назначения и сверяется средой на каждом вызове (I1).

\textbf{Digest контракта} отвечает на вопрос, что агент должен получить и по каким правилам результат принимается. Контракт содержит: родительскую цель; write scope; версии входных артефактов и зависимости; ожидаемые результаты и критерии приёмки; профиль подтверждаемых свойств: какое свойство требуется подтвердить, каким способом, для какого артефакта и при каких допущениях (например, сборка, регрессионные проверки и независимое ревью для одной задачи; машинно проверяемое доказательство свойства плюс ревью соответствия формулировки исходной задаче для другой); оценку в единицах сложности; класс риска; бюджеты. Контракт версионируется: любое изменение критериев или профиля создаёт новую версию и требует авторизации того же уровня, что и создание задачи. Обычная задача не может незаметно ослабить критерий собственного успеха.

\textbf{Адрес аттестации.} Финальное одобрение и каждая аттестация ссылаются на сочетание: версия контракта, кандидат $I$, версия политики проверки, digest окружения. Аттестация, ссылающаяся на устаревшую версию любого из четырёх, недействительна. Исход проверки есть одно из трёх: пройдена, провалена, не установлена; таймаут или сбой инструмента даёт не установлена и не является принятием.

\textbf{Иерархия.} Цель (подписана кор-нодой: область, класс риска, бюджет) содержит capability (выдана оркестратору: пути, типы изменений, лимиты), capability содержит задачу (контракт, план), задача содержит ревизию (диф внутри write scope). Каждый уровень только сужает предыдущий (I2); валидаторы проверяют вложенность на каждом переходе, среда проверяет digest плана на каждом вызове. Оркестратор декомпозирует полномочия, но не создаёт их.

\begin{center}
\begin{tikzpicture}[node distance=4mm]
\node[box,minimum width=56mm] (g) {Цель\\[1pt]\footnotesize подписана кор-нодой: область, класс риска, бюджет};
\node[box,minimum width=56mm,below=of g] (c) {Capability\\[1pt]\footnotesize выдана оркестратору: пути, типы изменений, лимиты};
\node[box,minimum width=56mm,below=of c] (t) {Задача\\[1pt]\footnotesize контракт (версия) и план (digest)};
\node[box,minimum width=56mm,below=of t] (r) {Ревизия и кандидат\\[1pt]\footnotesize диф внутри write scope; $I$ над головой $M$};
\draw[arr] (g) -- (c); \draw[arr] (c) -- (t); \draw[arr] (t) -- (r);
\node[lbl,right=8mm of c] {только сужение};
\end{tikzpicture}
\end{center}

\textbf{Классы риска.} Класс определяется детерминированно по путям, манифестам зависимостей и типам операций. Привилегированный класс (авторизация и доступ, зависимости, конфигурация сборки и CI, проверки, правила и манифест, подпись, публикация, секреты, миграции) требует двух независимых ревью либо одобрения человека из кор-команды, а его пути доступны для записи только из санбокса задачи этого класса. Семантически вредная задача внутри своей области остаётся остаточным риском; она ограничена областью, бюджетом, вестингом и остановкой.

\textbf{Write scope.} Две задачи с пересекающимися write scope не исполняются одновременно. Дизъюнктные области предотвращают конкурирующую запись в одни и те же пути, но не исключают семантических конфликтов между изменениями; их ловит только проверка кандидата (раздел 10). Область удерживается с момента назначения до интеграции или отмены задачи, так что следующая задача не начнёт менять те же пути, пока предыдущая ревизия ожидает принятия.

\textbf{Переиспользование.} Задача может использовать принятые артефакты других задач; зависимость записывается в журнал, и агент получает не просто полезный код, а его происхождение, контракт, зависимости и основания принятия.

\section{Правила валидаторов}

Детерминированные проверки на переходах процесса теперь задаёт конфигурация процесса: охранные условия ядра, инварианты модуля доставки и статический валидатор \texttt{curator playbook check} (curator-playbook v4.8, §6.3, §6.10, §10). Правила ниже остаются предикатами, которые кор-нода проверяет перед подписью записи журнала; где формулировки расходятся, права спецификация.

Кор-нода подписывает запись только после проверки детерминированных предикатов (I4):
\begin{enumerate}
\item переход состояния задачи допустим (раздел 10);
\item write scope лежит внутри capability, capability внутри цели; digest плана соответствует области и классу; версия контракта, на которую ссылается запись, есть текущая;
\item продьюсер, ревьювер и исполнитель проверок различны согласно политике независимости; если политика не выполнима, запись ожидания или эскалации, а не ослабление;
\item множество кандидатов в ревьюверы заморожено записью до появления хеша ревизии, случайность получена после, ревьювер вычислен из идентификатора задачи, хеша ревизии и случайности;
\item запись об интеграции содержит базу $M$, кандидата $I$, финальное одобрение и аттестации, адресованные ровно к $I$, текущей версии контракта, текущей политике и окружению, с исходом пройдена; текущая голова равна $M$;
\item назначение выдано действующей арендой с валидным fencing-токеном, в пределах лимитов capability и без пересечения write scope с удерживаемыми областями;
\item каждый хеш тела имеет сертификат доступности;
\item класс риска и требуемый им процесс соблюдены;
\item цель создана кор-нодой; предложения рабочих нод не исполняются до подписи кор-ноды;
\item операция раннера имеет идентификатор и ключ идемпотентности, а повтор ссылается на тот же идентификатор (раздел 11).
\end{enumerate}

Валидатор не проверяет качество кода и правильность ревью. Он проверяет, что ревью и проверки были, что они относятся ровно к тому состоянию и контракту, которые принимаются, и что их сделали независимые участники.

\section{Оркестратор}

Оркестратор планирует: декомпозирует цели в задачи с контрактами, назначает продьюсеров, инициирует ревью и интеграцию. Все его действия суть записи, проходящие правила раздела 8; у него нет доступа к среде помимо записей, он не выбирает ревьювера, не интегрирует в одиночку, не выходит за capability и не может заблокировать остановку.

\textbf{Аренда.} Роль оформлена записью аренды: нода, эпоха, случайный fencing-токен, предшествующая аренда, причина (создание, плановая передача, принудительный захват, восстановление), чекпоинт как база передачи, область. Аренда продлевается heartbeat; каждая запись оркестратора несёт эпоху и токен, валидаторы отвергают записи с устаревшей арендой и сохраняют их как расхождение (I7). Аренды с дизъюнктными областями могут действовать одновременно; задача на пересечении областей принадлежит аренде, покрывающей объединение, либо создаётся кор-нодой.

\textbf{Смена эпохи} происходит по таймауту heartbeat, исчерпанию ёмкости, решению кворума или нарушению правил; преемник определяется детерминированной ротацией. Плановая передача требует доказательства покоя (нет незавершённых вызовов, дочерних процессов и незавершённых операций раннера, дайджест хранилища стабилен), после чего сессия передаётся с записью хеша снимка. Принудительный захват есть отдельное явное действие: сессия не передаётся, предшественник отсекается токеном, преемник стартует из журнала. Преемник обязан считать переданную сессию недоверенной. Раз в $k$ эпох смена происходит холодно. Журнал содержит план, назначения и результаты, и этого достаточно для вступления: оркестратор пишет решения в журнал структурно, а снимки сессии есть чекпоинты с сертификатом доступности.

\textbf{Несколько операторов на одном борде.} Для борда процесса, который ведут несколько операторов, 28 сентября 2026 года решено направление: борд остаётся один на процесс, и у него или у совместной работы один принимающий писатель с fencing по эпохе. Каждая типизированная операция несёт актора, борд, версию конфигурации, область и эпоху аренды, ожидаемую ревизию и идентификатор операции. Эпоха проверяется при запуске агента, при финальной записи и при интеграции; писатель, потерявший полномочия, приостанавливает новые записи. Это та же схема, что аренда с fencing-токеном выше, применённая к одному борду. Решение идёт в трек waggle CM2; назначается ли принимающий писатель по владению процессом или выбирается участниками, пока открыто.

\section{Цикл задачи}

\begin{center}
\begin{tikzpicture}[node distance=4.5mm]
\node[box] (s1) {Авторизована};
\node[box,right=of s1] (s2) {Назначена};
\node[box,right=of s2] (s3) {Ревизия $R$};
\node[box,right=of s3] (s4) {Обратная связь};
\node[box,below=13mm of s1] (s5) {Кандидат $(M, I)$};
\node[box,right=of s5] (s6) {Ревью $I$};
\node[box,right=of s6] (s7) {Проверки $I$};
\node[box,right=of s7] (s8) {Интеграция};
\node[box,right=of s8] (s9) {Вестинг};
\draw[arr] (s1) -- (s2); \draw[arr] (s2) -- (s3); \draw[arr] (s3) -- (s4);
\draw[arr] (s4.south) -- ++(0,-5mm) -| (s5.north);
\draw[arr] (s5) -- (s6); \draw[arr] (s6) -- (s7); \draw[arr] (s7) -- (s8); \draw[arr] (s8) -- (s9);
\draw[arr] (s4.north) -- ++(0,5mm) -| node[lbl,above,pos=0.25] {не более $N$ раундов} (s3.north);
\draw[arr] (s7.south) -- ++(0,-4mm) -| node[lbl,below,pos=0.3] {отклонено или голова ушла вперёд: пересборка кандидата} (s5.south);
\end{tikzpicture}
\end{center}

\textbf{Авторизация и назначение.} Оркестратор декомпозирует цель в задачи с контрактами; задача совершается записью с QC. Рабочие ноды сообщают в heartbeat остаток ёмкости и окно сброса; оркестратор назначает продьюсера с учётом ёмкости, свободного слота и удерживаемых областей. Среда открывает санбокс по записи назначения.

\textbf{Ревизия и обратная связь.} Продьюсер работает в санбоксе и представляет хеш коммита $R$ над базой $B$, подписанный на своей машине (I6). Ранняя обратная связь ревьювера на $R$ допустима и полезна продьюсеру, но не является одобрением: одобрение относится только к кандидату. Число оплачиваемых раундов ограничено $N$; при превышении оркестратор выбирает из меню: передекомпозировать, сменить продьюсера, сменить ревьювера, эскалировать человеку, и выбор записывается.

\textbf{Кандидат.} Среда строит неизменяемого кандидата $I$ над текущей головой $M$ и записывает пару $(M, I)$. Множество кандидатов в ревьюверы замораживается до появления хеша $I$; случайность получается после, из порогового маяка кор-нод или проверяемой случайной функции [8]; ревьювер вычисляется из идентификатора задачи, хеша $I$ и случайности, и выбор пересчитывается любым валидатором.

\textbf{Ревью и проверки кандидата.} Ревьювер получает диф $M..I$ и подписывает одобрение, адресованное к $I$, версии контракта, политике и окружению. Проверки из профиля контракта выполняются на $I$ по инициативе ноды, третьей по отношению к продьюсеру и ревьюверу; критические наборы повторяет верификатор. Каждая аттестация имеет один из трёх исходов; принятие требует исхода пройдена по всем свойствам профиля.

\textbf{Интеграция.} Оркестратор вносит запись об интеграции с парой $(M, I)$; валидаторы проверяют правило 5; среда, получив QC, атомарно заменяет голову с $M$ на $I$. Если голова уже не равна $M$, сертификат неприменим: кандидат перестраивается, ревью и проверки повторяются (I3). Порядок задаёт очередь [14]. Перенос одобрения с $R$ на $I$ по совпадению patch-id не входит в базовую гарантию: одинаковый патч на другой базе может иметь другой смысл. Как расширение вне гарантии допустим перенос при детерминированном условии, что дельта базы $B..M$ не пересекается ни с write scope задачи, ни с множеством путей, которые продьюсер читал по receipts; проверки на $I$ повторяются всегда. Для моделей повторное ревью стоит токены, а не часы, поэтому базовая гарантия посильна там, где для людей она была бы неподъёмной.

\textbf{Вестинг и замена.} Начисление закрепляется, когда $I$ прожил $k$ эпох без отката. Замена продьюсера происходит при потере ноды, исчерпании ёмкости или по решению из меню; преемник получает сессию предшественника с правилом недоверия к ней.

\section{Внешние эффекты}

Операции раннера (подпись, публикация, устройства, данные) имеют эффекты вне среды, и журнал авторизаций сам по себе не определяет их семантику при сбоях. Поэтому:

\begin{itemize}
\item каждая операция имеет идентификатор и ключ идемпотентности; повтор ссылается на тот же идентификатор и параметры, и для получателей, поддерживающих идемпотентность, повтор не создаёт дополнительного эффекта;
\item исход операции записывается как выполнена, отклонена или не установлен; сбой раннера до записи исхода даёт не установлен;
\item при исходе не установлен слепой повтор запрещён: требуется сверка фактического состояния получателя либо разрешение человеком, и результат сверки записывается;
\item актуальность разрешения проверяется на границе исполнения: операция, авторизованная в прошлой эпохе или до остановки, не исполняется без новой авторизации; уже запущенная операция при смене аренды доводится до записи исхода, при остановке прерывается там, где это безопасно, иначе доводится и записывается.
\end{itemize}

\section{Рабочие ноды}

\textbf{Харнес владельца.} Собственного агентского цикла у системы нет. Рабочие ноды, как и оркестраторы, работают на существующих харнесах: Claude Code, Codex, Muse, харнес над локальной моделью. Независимость от харнеса есть принцип дизайна: владелец ёмкости сохраняет харнес, за который уже платит, и может сменить его, не теряя процессов, бордов и памяти (curator-keeper, \texttt{spec/keeper.md} v0.5.1, §7.7, решение 28 сентября 2026 года). Протокол ноды не зависит от того, какой цикл ведёт модель.

\textbf{Две формы клиента.} Публичному рою нужен донорский профиль ровно из двух привязок: удалённый shell среды и подпись ревизий; задача от проекта есть недоверенный ввод и не должна находить локальных инструментов для инжекта. Пока своего цикла нет, такой профиль держится переопределением инструментов существующего харнеса, которое curator-keeper §7.7 называет возможным первым путём к ограниченным рабочим агентам. Это слабее свойства конструкции: харнес по-прежнему несёт локальные инструменты, и профиль держится ровно настолько, насколько надёжно переопределение. Минимальный runtime, который не поставляет инструментов и делает донорский профиль свойством конструкции, отложен вместе с собственным циклом. Доверенная организация может подключить среду как сервер инструментов внутри родного харнеса вендора; донорский профиль там держится политикой организации. Обе формы говорят с проектом одним протоколом ноды; при присоединении нода предъявляет ключ и доказательства идентичности (раздел 13).

\textbf{Ёмкость любого происхождения.} Ресурс роя есть inference-ёмкость: подписки там, где условия вендора допускают, ключи API, локальные модели. Главный источник в публичном рое есть ёмкость, которая иначе сгорит: остаток окна квоты перед сбросом; heartbeat сообщает остаток, окно и класс стоимости, планировщик берёт истекающую ёмкость первой, нода ограничивает расход на проект за окно.

\textbf{Подпись и конфиденциальность.} Ключ ноды никогда не оказывается на среде; среда возвращает хеши, брокер на машине ноды подписывает ревизию и receipts (I6). Рабочая нода видит содержимое своего санбокса и отправляет его провайдеру своей модели; протокол не скрывает код от контрибьюторов и провайдеров, манифест объявляет допустимых провайдеров и условия хранения.

\section{Независимость и репутация}

\textbf{Три оси.} Независимость определяется по административному домену (против сговора и множественных личностей), семейству модели (против общих слепых зон) и тулчейну (против одинаковых ошибок среды). Проект объявляет политику: какие оси обязательны, какие сведения подтверждаются (ключ ноды и способ оплаты ёмкости подтверждают владельца в меру доверия к ним), какие принимаются на доверии (объявленное семейство модели), какие служат эвристикой. Различие строк в объявлении не доказывает различия владельцев, и политика должна это учитывать.

\textbf{Правило.} Отсутствие подходящих участников не ослабляет политику автоматически: задача ожидает, эскалируется либо получает отдельно авторизованное изменение условий, записанное в журнал. Независимость ревью и независимость исполнения проверок различаются: третья нода, запустившая проверки в той же среде, даёт разделение ролей, но не независимое подтверждение честности среды; его даёт только верификатор на другой машине.

\textbf{Идентичность и привязка ключей.} Ось владельца становится проверяемой через доказательства идентичности. Владелец ноды имеет долгоживущий ключ идентичности и доказывает его принадлежность внешней сущности двусторонне: заявлением, подписанным ключом, и артефактом на стороне провайдера с отпечатком ключа. Провайдеры суть резолверы с одним интерфейсом: по строке идентичности вернуть набор отпечатков и хеш свидетельства. Стартовый набор: TXT-запись домена под DNSSEC для организаций [20, 21], well-known по HTTPS, ключи аккаунта на хостинге кода; публикации в социальных платформах как необязательный класс; системы имён с блокчейном как плагины позже. Протоколу важен только результат резолвера и его свидетельство.

Резолв есть внешний вызов и недетерминирован между валидаторами, поэтому идентичность считается подтверждённой не потому, что кто-то проверил, а потому, что $t$ валидаторов записали одинаковый хеш свидетельства в одном окне; расхождение даёт не подтверждено. Свидетельство хранится в журнале с датой, и недоступность провайдера позже не отменяет зафиксированного наблюдения; каданс переподтверждения задаёт политика. Ключи нод короткоживущие и сертифицируются ключом идентичности записью в журнале: нода ротирует ключи, не переделывая доказательств, а репутация привязана к идентичности и переживает ключи и проекты. Отзыв есть удаление отпечатка у провайдера; первый валидатор, не нашедший его, записывает расхождение. Первое увиденное свидетельство пинится, и смена ключа в доказательстве есть событие, а не обновление.

Ревизии подписываются ключом ноды в формате SSH-подписи git; множество допустимых подписантов для среды выводится из журнала: ключ допустим, если сертифицирован подтверждённой идентичностью и назначен на задачу. Сообщения протокола уже подписаны ключами нод, добавляется только цепочка до ключа идентичности. Политика проекта назначает вес по классу доказательства: домен под DNSSEC выше публикации в социальной платформе; псевдонимная нода допускается с меньшим весом и большей глубиной ревью, но не запрещена. Симметрично для проектов: ключи кор-команды публикуются под доменом организации, и контрибьютор при присоединении проверяет подписантов манифеста через домен, а не на слово. Доказательство удостоверяет контроль над доменом или аккаунтом, а не честность владельца.

\textbf{Репутация.} Стоимость создания ноды равна стоимости ёмкости, что даёт естественный, не абсолютный барьер против множественных личностей [9]. После каждой задачи проект выдаёт ноде подписанную аттестацию исхода; репутация определяет глубину проверки, как адаптивная репликация в BOINC [6], и вес при выборе в ревьюверы. Пропорциональное вознаграждение при полной устойчивости к множественным личностям недостижимо [7], поэтому репутация снижает риск, а не устраняет его. Аттестации переносимы между проектами и взвешиваются репутацией выдавшего проекта.

\section{Остановка и федерация}

Любая кор-нода выпускает сертификат экстренной остановки: подпись одного ключа, эпоха и монотонный номер против повтора. Среда немедленно переходит в безопасный режим: не открывает санбоксы, не начинает операции раннера, не продвигает голову; позднее остановка упорядочивается журналом. Возобновление требует QC уровня. Злоупотребляющая нода исключается записью о составе. Откат журнала на версию до остановки невозможен (I8).

Несколько проектов образуют федерацию: контрибьютор сам выбирает, кому отдать ёмкость. Проект публикует подписанный манифест (адрес среды и хеш её демона, уровень доверия, формы клиента, ограничения на ёмкость и провайдеров, инвентарь, политика независимости), который компилируется в план контрибьютора при присоединении с показом разницы полномочий; подписанты манифеста проверяются через доказательства идентичности организации. Присоединение означает доверие среде и кор-команде проекта; обратное доверие не предполагается.

\section{Модель угроз}

\textbf{Участники и допущения.} Рабочая нода не доверена: злонамеренна, скомпрометирована или отравлена инжектом. Кор-нода и оркестратор: на L0 доверены; на L1 не более меньшинства скомпрометированных ключей при честном секвенсере; на L2 не более $f$ злонамеренных из $3f+1$. Среда доверена в отношении целостности, если не аттестована аппаратно. Провайдер модели видит содержимое санбокса и доверен в отношении конфиденциальности в пределах манифеста. Сеть не доверена; сообщения подписаны, повтор отсекается эпохой и номерами.

\textbf{Что гарантируется} при соблюдении допущений уровня: ни одного эффекта на среде без совершённой и актуальной авторизации (I1); ни одного продвижения головы без финального одобрения и аттестаций, адресованных ровно к продвигаемому кандидату и текущему контракту (I3); полномочия и условия приёмки задачи не превышают и не ослабляют выданного (I2); оркестратор не выбирает ревьювера и не интегрирует в одиночку; остановка доступна одной честной кор-ноде (I8); каждое обращение к среде через протокол имеет проверяемое происхождение (I10).

\textbf{Свидетельство и факт.} Receipt есть подписанное свидетельство среды, а не независимо установленный факт. Если среда скомпрометирована, её подпись не доказывает правдивость отчёта. Обнаружение компрометации по расхождению хешей условно: оно требует независимого наблюдения, верификатора на другой машине или собственных git-объектов у ревьювера, способных выявить расхождение.

\textbf{Чего не гарантируется.} Правильность суждения модели. Предотвращение компрометации среды. Конфиденциальность кода от контрибьюторов и провайдеров. Устойчивость к побегу из санбокса через уязвимость ядра. Единственность истории на L1 при злонамеренном секвенсере. Живучесть при злонамеренном большинстве кор-нод. Полнота наблюдения за недоверенной нодой. Честность владельца подтверждённой идентичности: доказательство удостоверяет контроль над доменом или аккаунтом, а не поведение. Справедливость экономики.

\section{Свойства и оценки}

\textbf{Целостность.} На L0 целостность сводится к одному ключу и среде. На L1 запись требует $t > n/2$ подписей: компрометация меньшинства ключей не совершает записи, а единственность истории обеспечивает секвенсер. На L2 журнал сохраняет целостность и единственность истории, пока злонамеренных кор-нод не более $f$ [1, 2]. На всех уровнях среда не исполняет ничего без QC, а откат исключён высшей принятой версией.

\textbf{Живучесть.} Секвенсер даёт продвижение журнала, пока доступны он и достаточное для QC число валидаторов. Остановка есть намеренная потеря живучести в пользу целостности. Политика независимости и удержание областей тоже покупают целостность ценой ожидания, и это ожидание измеряется (раздел 18).

\textbf{Рабочий слой.} Пусть доля злонамеренных рабочих нод $p$, честный ревьювер обнаруживает злонамеренного кандидата с вероятностью $d$, ревьюверов $k$. Кандидат проходит ревью с вероятностью $P = q^k$, где $q = p + (1-p)(1-d)$: при $p = 0{,}1$ и $d = 0{,}7$ получаем $0{,}37$, $0{,}14$, $0{,}05$ для $k = 1, 2, 3$; при $d = 0{,}9$ получаем $0{,}19$, $0{,}04$, $0{,}01$. Формула предполагает независимость ошибок, которую три оси делают лишь менее нарушенной; это иллюстрация того, что $d$ влияет сильнее $p$, а не оценка безопасности.

\textbf{Измерение $d$.} Вероятность $d$ больше не обязана оставаться допущением. Калибровка ревьюверов посевом дефектов (curator-playbook, \texttt{spec/process-configuration.md} v4.8, §6.16, предложено) даёт ревьюверу копии кандидатов с запечатанными посеянными дефектами известных классов и считает долю найденных; обозначим её $r$. Доля хранится для роли на конкретной модели, с её усилием и профилем роли, по классам и в целом, как измеряемый термин потребности \texttt{recall} (§9 той же спецификации). Для честного ревьювера $r$ есть оценка $d$, а $1 - r$ есть оценка того, что он пропустит посаженную уязвимость. Тогда $q$ и $P = q^k$ перестают быть иллюстрацией и становятся эмпирической оценкой для заданных ролей и моделей при известной доле $p$. Оговорка та же, что в спецификации: оценка верна настолько, насколько посеянные дефекты так же трудно найти, как настоящие, а уязвимость, которую злонамеренный продьюсер прячет намеренно, может оказаться труднее любого посева; классы, взятые из дефектов, прошедших прежние ревью, сближают их. Калибровка измеряет $d$, но не $p$ и не совместный пропуск: допущение независимости остаётся, пока $k$ ревьюверов не калибруются на одних и тех же посевах, где совместный пропуск можно считать прямо.

\section{Проверяемый фрагмент}

Инварианты и правила валидаторов детерминированы, поэтому ключевой фрагмент протокола допускает исполняемую модель. Мы построили модель фрагмента кандидат, аттестации, запись об интеграции, атомарное продвижение головы и проверили её исчерпывающим перебором состояний в заданных границах.

\textbf{Модель.} Две задачи: $A$ с полным поведением (контракт может менять версию, записи могут нести устаревшую эпоху) и $B$, которая только интегрируется и тем сдвигает голову. Состояние: голова, эпоха аренды, и для каждой задачи наличие ревизии, база кандидата, аттестация (адрес и версия контракта), версия контракта, признак интеграции и запись об интеграции с числом доставок. Действия: представить ревизию; построить кандидата над текущей головой; аттестовать кандидата; изменить версию контракта; сменить эпоху; предложить запись об интеграции с текущей или устаревшей эпохой; доставить запись, в том числе повторно. Границы: версии контракта не более двух, эпохи не более двух, каждая запись доставляется не более двух раз. Правила при доставке: (cas) голова равна базе записи; (epoch) запись несёт текущую эпоху; (contract) аттестация ссылается на текущую версию контракта; (bind) аттестация адресована ровно к кандидату. Свойство, проверяемое при каждом продвижении головы:

\begin{quote}
голова продвигается к кандидату, построенному над головой в этот момент; аттестация адресована к этому кандидату и ссылается на текущую версию контракта; запись несёт текущую эпоху; ни одна запись не применяется дважды.
\end{quote}

\textbf{Результаты.} При всех четырёх правилах достижимо 1090 состояний, и свойство выполнено во всех. При отключении любого одного правила перебор находит контрпример:

\begin{center}
\small
\begin{tabular}{p{2.0cm} p{6.4cm} p{4.6cm}}
\toprule
Без правила & Нарушение & Кратчайшая трасса \\
\midrule
cas & одна запись применена дважды при повторной доставке; и продвижение к кандидату, построенному не над текущей головой, после интеграции другой задачи & 6 действий; 11 действий \\
epoch & запись оркестратора устаревшей эпохи продвинула голову & 6 действий \\
contract & аттестация под прежней версией контракта принята после изменения критериев & 6 действий \\
bind & аттестация, сделанная на базе до сдвига головы, перенесена на кандидата над новой головой & 11 действий \\
\bottomrule
\end{tabular}
\end{center}

Последняя строка есть ровно сценарий переноса одобрения по patch-id: одинаковое изменение, другая база, аттестация без адреса кандидата. Модель показывает, что правило привязки к кандидату не оптимизация, а условие свойства.

\textbf{Границы результата.} Это проверка ограниченного пространства сценариев, а не доказательство для всех допустимых исполнений: границы по числу задач, версий и доставок заданы явно, среда и валидаторы в модели честны, а суждение и содержимое кода не моделируются. Трассы модели используются как векторы соответствия для демона среды и секвенсера; следующий шаг есть модель в языке спецификаций с механической проверкой [17] и распространение на аренду и остановку. Модель прилагается к работе.

\section{Эксперименты}

Мы разделяем две проверки, потому что они отвечают на разные вопросы.

\textbf{Соблюдение протокола.} Реализация подвергается сбоям и нарушителям, воспроизводимым программно: рабочая нода, выходящая за write scope; запись оркестратора с устаревшей арендой; подмена кандидата после аттестации; изменение критериев после ревью; повторная доставка записи об интеграции; операция раннера с потерянным исходом; смерть оркестратора с удалением его сессии. Критерий: ни одно нарушение не приводит к эффекту на среде или продвижению головы, а восстановление проходит по журналу. Качество модели здесь несущественно.

\textbf{Полезность роя.} Единица сравнения есть заранее определённая внешняя задача или завершённая функция проекта с одинаковыми критериями приёмки, а не внутренние задачи оркестратора, число которых он сам может менять. Три условия при сравнимых ресурсных ограничениях: одиночный агент; параллельный процесс без протокола (агенты одного владельца с общим репозиторием и блокировками); развёртывание L0. Стоимость включает inference, исполнение проверок, занятость среды и человеческое вмешательство. Отдельно измеряются ожидание слотов, блокировки областей, повторная работа после смены головы, восстановление после потери агента, доля пересмотренных отклонений и доля откатов через $k$ эпох. Ориентир масштаба задаёт эксперимент, в котором три личные подписки за три дня совместно формализовали крупный результат [16]; первое число на коде должно быть получено в том же масштабе.

Отдельный эксперимент проверяет I7 непосредственно: остановить оркестратора, удалить его сессию и проверить, продолжает ли преемник работу из журнала и доступных артефактов.

\section{Смежные работы}

Prove2Me [15] есть открытая платформа коллаборативной формализации: пользователи запускают миссии, агенты любых владельцев вносят доказательства, граф формулировок задаёт фронт работы, формулировки отделены от доказательств, описание каждого узла делает результаты находимыми и переиспользуемыми. На ней показана совместная формализация крупных результатов, в том числе тремя личными подписками за три дня [16]. Это ближайшая работа по объединению чужих агентов, и от неё нас отличает домен: там принятие есть проверка ядром, у нас его приходится конструировать; наш контракт задачи и граф зависимостей есть перенос разделения формулировок и доказательств на код. Эксперимент с командой параллельных агентов, построившей компилятор [18], есть ближайшая работа по обычной разработке: агенты в отдельных контейнерах, общий репозиторий, блокировки задач, самостоятельная интеграция; там все агенты одного владельца и доверены, у нас нет. Византийски устойчивые протоколы [1, 2, 4, 5] и Raft [3] задают кор-слой и модель снимка и лога; очереди интеграции [14] правило никогда не сливать непроверенное; API удалённого исполнения сборок [19] модель действия, адресуемого по содержимому; журналы прозрачности [11, 12] и дерево Меркла [10] форму журнала; аттестации цепочки поставки [13] принцип, кто и на каком шаге вправе произвести артефакт; BOINC [6] модель добровольного вклада с адаптивной репликацией. Публикация ключей через DNS [20, 21] и системы доказательств идентичности, в которых отпечаток ключа размещается у стороннего провайдера, а владелец подписывает обратное заявление, задают форму привязки ключей нод к внешним идентичностям.

\section{Заключение}

Мы описали протокол совместного изменения проекта агентами независимых владельцев в домене без механического оракула принятия, и сформулировали его как набор точных условий, механизмов и проверок. Среда одна, ключи у владельцев, полномочия только в журнале и только сужаются, условия приёмки версионированы, голова продвигается ровно к тому кандидату, к которому адресованы одобрение и аттестации, суждение вынесено за консенсус и сделано независимым и подотчётным. Для каждого уровня доверия сказано, от чего защищает сертификат. Фрагмент интеграции проверен исчерпывающим перебором в заданных границах, и модель показывает, что каждое из четырёх правил необходимо. Остальное подлежит эксперименту, который может нас опровергнуть. Примитивы для реализации существуют; работа состоит в их безопасной композиции.

\section*{Приложение. Учёт вклада}

Мера вклада фиксируется до работы: оценка задачи при авторизации контракта. Задача оплачивается за принятый результат, а не за токены. Оценка делится на пулы продьюсеров, ревью и оркестрации по правилам, вычислимым из записей: пул продьюсеров получает продьюсер принятой ревизии, заменённый получает часть, если принятый кандидат построен на его коммитах; раунды сверх $N$ и отклонения, проигранные в споре, не оплачиваются; доли финального одобрившего и оркестратора вестятся вместе с задачей. Ретроактивная корректировка кворумом ограничена заранее оговорённой долей. Во что конвертируются единицы вклада, определяет проект. Стимулы ревью есть открытый вопрос; экономика вынесена в отдельную работу.

\section*{Литература}
\begin{english}
\begin{enumerate}[label={[\arabic*]},leftmargin=2.6em,itemsep=2pt]
\item L. Lamport, R. Shostak, M. Pease. The Byzantine Generals Problem. ACM Transactions on Programming Languages and Systems, 4(3), 1982.
\item M. Castro, B. Liskov. Practical Byzantine Fault Tolerance. Proceedings of OSDI, 1999.
\item D. Ongaro, J. Ousterhout. In Search of an Understandable Consensus Algorithm. Proceedings of USENIX ATC, 2014.
\item E. Buchman. Tendermint: Byzantine Fault Tolerance in the Age of Blockchains. M.Sc. thesis, University of Guelph, 2016.
\item M. Yin, D. Malkhi, M. K. Reiter, G. Golan Gueta, I. Abraham. HotStuff: BFT Consensus with Linearity and Responsiveness. Proceedings of PODC, 2019.
\item D. P. Anderson. BOINC: A System for Public-Resource Computing and Storage. Proceedings of the 5th IEEE/ACM International Workshop on Grid Computing, 2004.
\item A. Cheng, E. Friedman. Sybilproof Reputation Mechanisms. Proceedings of the ACM SIGCOMM Workshop on Economics of Peer-to-Peer Systems, 2005.
\item S. Micali, M. Rabin, S. Vadhan. Verifiable Random Functions. Proceedings of FOCS, 1999.
\item J. R. Douceur. The Sybil Attack. Proceedings of IPTPS, 2002.
\item R. C. Merkle. A Digital Signature Based on a Conventional Encryption Function. Advances in Cryptology, CRYPTO 1987.
\item B. Laurie, A. Langley, E. Kasper. Certificate Transparency. RFC 6962, 2013.
\item Z. Newman, J. S. Meyers, S. Torres-Arias. Sigstore: Software Signing for Everybody. Proceedings of ACM CCS, 2022.
\item S. Torres-Arias, H. Afzali, T. K. Kuppusamy, R. Curtmola, J. Cappos. in-toto: Providing Farm-to-Table Guarantees for Bits and Bytes. Proceedings of USENIX Security, 2019.
\item G. Hoare. The Not Rocket Science Rule of Software Engineering. 2014.
\item S. Chen et al. Prove2Me: An Open Collaborative Platform for Scaling Math Formalization. arXiv:2608.28433, 2026.
\item Anthropic. Formalizing Fermat's Last Theorem. anthropic.com/research, 2026.
\item L. Lamport. Specifying Systems: The TLA+ Language and Tools for Hardware and Software Engineers. Addison-Wesley, 2002.
\item N. Carlini. Building a C Compiler with a Team of Parallel Claudes. Anthropic Engineering, 2026.
\item Remote Execution API. github.com/bazelbuild/remote-apis.
\item P. Hoffman, J. Schlyter. The DNS-Based Authentication of Named Entities (DANE) Transport Layer Security (TLS) Protocol: TLSA. RFC 6698, 2012.
\item P. Wouters. DNS-Based Authentication of Named Entities (DANE) Bindings for OpenPGP. RFC 7929, 2016.
\end{enumerate}
\end{english}

\end{document}
